\documentclass{llncs}

\usepackage[utf8]{inputenc}
\usepackage[T1]{fontenc}
\usepackage[brazilian]{babel}
\usepackage{csquotes}
\usepackage{hyphenat}
\hyphenation{pro-ble-ma}
\usepackage{imakeidx}
\makeindex[columns=3, title=Alphabetical Index, intoc]
\usepackage{enumerate}
\usepackage{amssymb}
\usepackage{amsmath}
\usepackage{array}

% Better handle of numeric things.
% Configured to French style (same 
% as Brazilian conventions)
\usepackage[locale=FR]{siunitx}
\usepackage{hyperref}
\usepackage{graphicx}
\graphicspath{ {images/} }
\usepackage{longtable}
\usepackage{listings}
\setcounter{tocdepth}{1}
\usepackage{tocbibind}

\begin{document}
\begin{titlepage}           
\end{titlepage}
%%%%%%%%%%%%%%%%%%% TITLE PAGE
\begin{titlepage}
  \begin{center}
    \vspace*{1cm}
    %%% TITLE
    \Huge
    \textbf{Tarefa Prática 2}
    \vspace{0.5cm}
    
    %%% CURRICULAR UNIT
    \LARGE
    Modelação e Caracterização de Tráfego
    %%% AUTHORS    
    \vspace{1.0cm}
    \small
    \textbf{\\PG39254 - Igor Araújo\\PG39255 - Matheus Gonçalves\\PG41017 - I-Ping}
    
    %%% LOGO
    \vspace{1.0cm}
    \begin{figure}[ht]
    \includegraphics[width=0.8\textwidth]{uminho.jpg}
    \centering
    \end{figure}
    
    % FOOTER
    \vspace{4.5cm}
    Departamento de Informática\\
    Universidade do Minho\\
    Braga - Portugal\\
    \today
          
  \end{center}
\end{titlepage}

\tableofcontents

\clearpage

\section{Introdução}

%
Insira sua Introdução aqui  \cite{castro}. Insira sua Introdução aqui  Insira sua Introdução aqui   Insira sua Introdução aqui \cite{racket}   Insira sua Introdução aqui   Insira sua Introdução aqui   Insira sua Introdução aqui \cite{deb:agra}   Insira sua Introdução aqui  

%
\section{Metodologia}
%
Aqui está uma metodologia \emph{incrível}. Aqui está uma metodologia  
Aqui está uma metodologia Aqui está uma metodologia Aqui está uma metodologia Aqui está uma metodologia 
Aqui está uma metodologia 
%

\subsection{Os Modelos}
%
Os modelos são caracteres \enquote{X} representados em grades quadradas em três tamanhos diferentes, como mostrado na figura~\ref{fig:xs}.

\subsection{Testes}
Testes Testes Testes Testes Testes Testes Testes Testes Testes Testes Testes Testes Testes Testes Testes Testes Testes Testes \par Testes Testes Testes Testes Testes Testes Testes Testes Testes Testes Testes Testes Testes Testes
\begin{enumerate}
  \item test1
  \item test2
  \item test3
\end{enumerate}
\renewcommand{\labelenumiii}{\roman{enumiii}}
\begin{enumerate}[i]
  \item test1
  \item test2
  \item test3
\end{enumerate}

\begin{figure}[ht]
\includegraphics[scale=0.5]{uminho.jpg}
\centering
\caption{\enquote{X} em modelos de tamanhos diferentes}
\label{fig:xs}
\end{figure}

%
\section{Resultados}
%
Os resultados dos experimentos se encontram na tabela 1. 

\begin{table}
\centering
\begin{tabular}{|c|S|c|c|c|c|c|c|c|c|}
\hline
\multicolumn{3}{|c|}{Parâmetros} & \multicolumn{5}{|c|}{Geração até chegar à solução} & \multicolumn{2}{|c|}{Desempenho} \\ 
\hline
Grade & Mutação & População & \multicolumn{2}{|c|}{95\% de confiança} & 1º Quartil & Mediana & 3º Quartil & Indivíduos & IC \\ 
\hline
3x3 & 0.1 & 10 & 2 & 1000 & 16 & 167 & 435 & 1670 & 96,17\%\\ 
\hline
3x3 & 0.01 & 10 & 1 & 1000 & 12 & 132 & 383 & 1320 & 92,42\%\\ 
\hline
3x3 & 0.001 & 10 & 2 & 1000 & 40 & 197 & 430 & 1970 & 97,87\%\\ 
\hline
3x3 & 0 & 10 & 2 & 1000 & 32 & 135 & 404 & 1350 & 92,85\%\\ 
\hline
3x3 & 0.1 & 100 & 1 & 6 & 2 & 3 & 4 & 300 & 44,37\%\\ 
\hline
3x3 & 0.01 & 100 & 1 & 7 & 2 & 3 & 4 & 300 & 44,37\%\\ 
\hline
3x3 & 0.001 & 100 & 1 & 7 & 2 & 3 & 4 & 300 & 44,37\%\\ 
\hline
3x3 & 0 & 100 & 1 & 8 & 2 & 3 & 4 & 300 & 44,37\%\\ 
\hline
3x3 & 0.1 & 1000 & 1 & 2 & 1 & 1 & 1 & 1000 & 85,84\%\\ 
\hline
3x3 & 0.01 & 1000 & 1 & 2 & 1 & 1 & 1 & 1000 & 85,84\%\\ 
\hline
3x3 & 0.001 & 1000 & 1 & 2 & 1 & 1 & 1 & 1000 & 85,84\%\\ 
\hline
3x3 & 0 & 1000 & 1 & 3 & 1 & 1 & 1 & 1000 & 85,84\%\\ 
\hline
4x4 & 0.1 & 10 & 406 & 1000 & 1000 & 1000 & 1000 & 10000 & \\ 
\hline
4x4 & 0.01 & 10 & 535 & 1000 & 1000 & 1000 & 1000 & 10000 & \\ 
\hline
4x4 & 0.001 & 10 & 241 & 1000 & 1000 & 1000 & 1000 & 10000 & \\ 
\hline
4x4 & 0 & 10 & 185 & 1000 & 1000 & 1000 & 1000 & 10000 & \\ 
\hline
4x4 & 0.1 & 100 & 5 & 27 & 11 & 14 & 17 & 1400 & 2,11\%\\ 
\hline

\hline
\hline\end{tabular}
\label{tab:resultados}
\caption{\small{Resultados brutos}} 
\end{table}

\newpage
\section{Anexo I}
\label{sec:anexo_I}

\begin{table}
\centering
\begin{longtable}{p{.20\textwidth} p{.20\textwidth} p{.20\textwidth} p{.40\textwidth}}
\label{tbl:2}

%\begin{tabular}{p{2cm}|p{2cm}|p{2cm}|p{8cm}}

\\\textbf{Test} & \textbf{Metric} & \textbf{Plataform} & \textbf{Description}\\
\hline\hline
Download (TCP) & Download speed & Whiteboxes, Routers, Android, iOS & The download speed in Mbps when downloading (using TCP) random bytes from a test server \\ 
\hline 
  & TCP Retransmissions & Whiteboxes, Routers & The number of retransmitted TCP segments/packets \\ 
\hline
  & Burst download speed & Whiteboxes, Routers & The download speed during the first 5 seconds of a test \\ 
\hline 
  & Sustained download speed & Whiteboxes, Routers & The download speed of the test during the last 5 seconds \\ 
\hline 
  & Percentage of Best & Whiteboxes, Routers & Download speed result as a percentage of the user's best ever result \\ 
\hline 
  & Percentage of Advertised & Whiteboxes, Routers & Download speed result as a percentage of their package's advertised downstream speed \\ 
\hline 
Download (HTML5) & Download speed & Web & The download speed in Mbps when downloading (using TCP) random bytes from a test server using HTML5 APIs(WebSockets and Fetch) \\ 
\hline 
Download (Lightweight UDP) & Download speed & Whiteboxes, Routers & The download speed in Mbps when downloading (using UDP) from a test server, using less data than the TCP test \\ 
\hline 
Download (Hardware acceleratedUDP) & Download speed & Broadcom-based Routers & The download speed in Mbps when downloading (using UDP) random bytes from a test server \\
\hline 
\caption{Tabela com alguns exemplos}
\end{longtable}
\end{table}
%
\section{Conclusões}
%
O experimento é pequeno para mostrar dados conclusivos, mas mostra indícios do comportamento dos parâmetros de maneira bastante consistente. Obviamente, alterações no algoritmo ou alterações na forma como o problema é representado devem alterar esse comportamento.

Futuros trabalhos podem ser feitos usando-se uma metodologia parecida, mas com modificações no algoritmo e no problema para validar os dados aqui obtidos.

%
% ---- Bibliography ----
%
\begin{thebibliography}{5}
%

\bibitem {castro}
de Castro, L.N.: Fundamentals of Natural Computing: Basic Concepts, Algorithms, and Applications. CRC Press (2006).

\bibitem{racket}
Felleisen, M., Findler, R.B., Flatt, M.: The Racket Manifesto. LIPIcs-Leibniz. (2015).

\bibitem{deb:agra}
Deb, K., Agrawal, S.: Understanding interactions among genetic algorithm parameters. Foundations of Genetic Algorithms. (1999).

\end{thebibliography}

\end{document}
